\documentclass[12pt,a4paper]{report}

% --- Paquetes de Idioma y Codificación ---
\usepackage[utf8]{inputenc}
\usepackage[T1]{fontenc}
\usepackage[spanish, es-tabla]{babel}

% --- Diseño de Página ---
\usepackage{geometry}
\geometry{top=2.5cm, bottom=2.5cm, left=3cm, right=3cm}

% --- Matemáticas y Símbolos ---
\usepackage{amsmath, amssymb, amsthm}

% --- Colores y Cajas ---
\usepackage[dvipsnames]{xcolor}
\usepackage[most]{tcolorbox}

% Caption

\usepackage{caption}

% Gráficos

\usepackage{graphicx}

% Espaciadora

\newcommand{\esp}{\vspace{0.5cm}}

% promedio

\newcommand{\prom}[1]{\left\langle #1 \right\rangle}

% ket

\newcommand{\ket}[1]{| #1 \rangle}

% bra

\newcommand{\bra}[1]{\langle #1 |}

% braket

\newcommand{\braket}[2]{\langle #1 | #2 \rangle}

% Bibliografía

\usepackage[style=apa, backend=biber]{biblatex} % Estilo APA
\addbibresource{biblio.bib} % Nombre de tu archivo .bib

% Configuración de cajas personalizadas
\newtcolorbox{importante}{
	colback=red!5!white,
	colframe=red!75!black,
	title=Importante,
	fonttitle=\bfseries
}

\newtcolorbox{nota}{
	colback=blue!5!white,
	colframe=blue!75!black,
	title=Nota,
	fonttitle=\bfseries
}

% Definición de la caja de bibliografía
\newtcolorbox{bibliografia}{
	colback=gray!5!white,
	colframe=gray!75!black,
	title=Bibliografías útiles,
	fonttitle=\bfseries,
	breakable
}

% Socrative para óptica 2

\newtcolorbox{Socrative}{
	colback=orange!5!white,
	colframe=orange!75!black,
	title=Socrative,
	fonttitle=\bfseries,
	breakable
}

% Float

\usepackage{float}

% --- Hipervínculos ---
\usepackage{hyperref}
\hypersetup{
	colorlinks=true,
	linkcolor=blue,
	filecolor=magenta,      
	urlcolor=cyan,
}

% --- Datos del Documento ---
\title{Óptica II}
\author{Chengyu Jin}
\date{\today}

\begin{document}
	
	\maketitle
	
	\tableofcontents
	\newpage
	
	\chapter{Quantum Ideal Gas}
	
	\section{Non-interacting Ising model}
	
	\noindent $H = -J \sum_{i} s_i$ where $s_i = \pm 1$ (spin)
	
	\esp
	
	\noindent We must to find:
	
	\esp
	
	\noindent $Z$ \\
	\noindent $A$ \\
	\noindent $U = \langle H \rangle$ \\
	\noindent $\langle M \rangle = \langle \sum s_i \rangle / N$
	
	\esp
	
	\noindent \textbf{For interacting case}
	
	\esp
	
	At $T_c$ my particles are disordered. Below $T_c$, symmetry is broken (spontaneous magnetization).
	
	\esp
	
	$Z = (z_i)^N$ 
	
	\esp
	
	$z_i = \sum_{s_i = \pm 1} e^{-\beta H_i} = e^{\beta J} + e^{-\beta J} = 2\cosh(\beta J)$
	
	\esp
	
	$Z = (2\cosh(\beta J))^N$
	
	\esp
	
	$A = -kT \ln Z = -NkT \ln(2\cosh(\beta J))$
	
	\esp
	
	$U = \left( -\frac{\partial \ln Z}{\partial \beta} \right)_{N, V} = -\frac{\partial}{\partial \beta} (N \ln(2\cosh(\beta J))) = -N \frac{2\sinh(\beta J)}{2\cosh(\beta J)} \cdot J = -NJ \tanh(\beta J)$
	
	\esp
	
	\noindent \textbf{In other way:}
	
	\esp
	
	$U = \langle H \rangle = -J \sum_i \langle s_i \rangle = -NJ \tanh(\beta J)$ \\
	$\langle s_i \rangle = 1 \cdot P(s_i=1) - 1 \cdot P(s_i=-1) = \frac{2\sinh(\beta J)}{2\cosh(\beta J)} = \tanh(\beta J)$
	
	\esp
	
	where $P(s_i) = \frac{e^{-\beta J s_i}}{2\cosh(\beta J)}$ \\
	$P(1) = \frac{e^{\beta J}}{2\cosh(\beta J)}$ \\
	$P(-1) = \frac{e^{-\beta J}}{2\cosh(\beta J)}$
	
	\section{Ideal Quantum Gas}
	
	System made of $N$ identical particles:
	\begin{equation}
		\Psi(r_1, r_2, \dots, r_N, t)
	\end{equation}
	
	\esp
	
	Its evolution is given by:
	\begin{equation}
		i\hbar \frac{\partial \Psi}{\partial t} = H\Psi \quad ; \quad H = \sum_{i=1}^N \frac{p_i^2}{2m} + V(r_1, r_2, \dots, r_N)
	\end{equation}
	
	\esp
	
	$P_{ij}$ is the \textbf{permutation operator} which exchanges the labels of particles $i$ and $j$:
	\begin{equation}
		P_{ij} \Psi(\dots, r_i, \dots, r_j, \dots, t) = \Psi(\dots, r_j, \dots, r_i, \dots, t)
	\end{equation}
	
	\esp
	
	$[H, P_{ij}] = 0 \quad \forall i, j \quad (i \neq j)$
	
	\esp
	
	If $H\phi_n = E_n\phi_n$:
	\begin{equation}
		\phi(r_1, \dots, r_N, t) = \sum a_n \phi_n(r_1, \dots, r_N) \quad ; \quad a_n = \langle \phi_n | \Psi \rangle
	\end{equation}
	
	\esp
	
	Now we apply $P_{ij}$ to $\Psi$ and then $H$:
	\begin{align*}
		H P_{ij} \Psi(\dots, r_i, \dots) &\to \\
		H \Psi(\dots, r_j, \dots, r_i, \dots, t) &= \sum a_n E_n \phi_n(\dots, r_j, \dots, r_i, \dots, t) = \sum a_n E_n \underbrace{P_{ij}}_{\text{scalars}} \phi_n(\dots, r_i, \dots, r_j, \dots, t) \\
		&= P_{ij} \sum a_n E_n \phi_n(\dots, r_i, \dots, r_j, \dots, t) = P_{ij} H \Psi \\
		&\to \text{We have proved that } [H, P_{ij}] = 0 \quad \forall i, j \quad (i \neq j)
	\end{align*}
	
	\esp
	
	Commuting observables have a complete set of common eigenfunctions:
	\begin{equation}
		P_{ij} H \phi = P_{ij} E \phi = E (P_{ij} \phi) \xrightarrow{[H, P_{ij}]=0} H (P_{ij} \phi)
	\end{equation}
	This implies $P_{ij} \phi$ is also an eigenfunction of $H$ with the eigenvalue $E$.
	
	\esp
	
	$P_{ij} \phi = \Theta \phi$
	
	\esp
	
	\textbf{Value of $\Theta$?}
	\begin{align*}
		P_{ij} \phi(\dots, r_i, \dots, r_j, \dots, t) &= \phi(\dots, r_j, \dots, r_i, \dots, t) \\
		P_{ij}^2 \phi(\dots, r_i, \dots, r_j, \dots, t) &= \phi(\dots, r_i, \dots, r_j, \dots, t) = I \phi
	\end{align*}
	Then $P_{ij}^2 = I$.
	
	\esp
	
	As $P_{ij} \phi = \Theta \phi$:
	\begin{align*}
		P_{ij} \phi(\dots, r_i, \dots, r_j, \dots, t) &= \Theta \phi(\dots, r_i, \dots, r_j, \dots, t) \\
		P_{ij}^2 \phi(\dots, r_i, \dots, r_j, \dots, t) &= \Theta^2 \phi(\dots, r_i, \dots, r_j, \dots, t)
	\end{align*}
	As $P_{ij}^2 = I$, then $\Theta^2 = 1 \implies \Theta = \pm 1$.
	
	\esp
	
	\begin{itemize}
		\item \textbf{Bosons:} $\Theta = +1 \to$ \textit{Symmetric}
		\item \textbf{Fermions:} $\Theta = -1 \to$ \textit{Antisymmetric}
	\end{itemize}
	
	\esp
	
	\textbf{What happens if two particles are in the same position ($r_i = r_j$) and I exchange them with $P_{ij}$?}
	
	\esp
	
	\textbf{First case (Fermions):} $P_{ij} \phi(\dots, r_i, \dots, r_i, \dots, t) = -\phi(\dots, r_i, \dots, r_i, \dots, t)$ \\
	$\implies \phi(\dots, r_i, \dots, r_i, \dots, t) = 0 \quad \to$ It is \textbf{forbidden} that two particles occupy the same position. These particles are known as \textbf{Fermions} (half-integer spin).
	
	\esp
	
	\textbf{Second case (Bosons):} $P_{ij} \phi(\dots, r_i, \dots, r_i, \dots, t) = \phi(\dots, r_i, \dots, r_i, \dots, t)$ \\
	$\implies \phi(\dots, r_i, \dots, r_i, \dots, t) = \phi(\dots, r_i, \dots, r_i, \dots, t) \quad \to$ It is \textbf{not forbidden} that two particles occupy the same position. These particles are known as \textbf{Bosons} (integer spin).
	
	\esp
	
	\textbf{Examples:}
	
	\esp
	
	\begin{table}[h]
		\centering
		\begin{tabular}{ll}
			\textbf{Fermions} & \textbf{Bosons} \\
			\textcolor{red}{Electron} & \textcolor{red}{Quark} \\
			\textcolor{red}{Proton} & \textcolor{red}{Gluon} \\
			\textcolor{red}{Neutron} & \textcolor{red}{$\alpha$ particle} \\
			$\vdots$ & $\vdots$ \\
			$\phi$ is antisymmetric ($\Theta = -1$) & $\phi$ is symmetric ($\Theta = +1$)
		\end{tabular}
	\end{table}
	
	\esp
	
	\textbf{Ideal $\to H = \sum_i^N H_i$}
	
	\esp
	
	$\sum_i n_i = N$ \\
	$\sum_i n_i \epsilon_i = E$
	
	\esp
	
	\textbf{Fermions vs Bosons occupancy:}
	
	\esp
	
	\begin{itemize}
		\item \textbf{Fermions:} $n_j \in \{0, 1\}$ (Pauli exclusion principle).
		\item \textbf{Bosons:} $n_j \in \{0, 1, 2, \dots, N\}$ (No restriction).
	\end{itemize}
	
	\esp
	
	\begin{equation}
		Z(T, V) = \sum_{\{n_i\}} e^{-\beta \sum n_i \epsilon_i} = \sum_{\{n_i\}} e^{-\beta E(\{n_i\})} \delta_{N, \sum n_i}
	\end{equation}
	$\to$ \textit{This restriction ($\sum n_i = N$) does not allow us to perform the sum easily.}
	
	\esp
	
	\textbf{Boltzmann approximation} $\to z_N = \sum_{\{n_i\}} \frac{n_1! n_2! \dots n_i!}{N!} e^{-\beta \sum n_i \epsilon_i} \approx \frac{1}{N!} \sum_{\{n_i\}} e^{-\beta \sum n_i \epsilon_i}$
	
	\esp
	
	\textbf{Ensembles Comparison:}
	
	\esp
	
	\begin{itemize}
		\item \textbf{Microcanonical:} $(E, V, N)$ fixed. Isolated system.
		\item \textbf{Canonical:} $(\langle E \rangle, V, N)$ fixed. Thermal contact with reservoir at $T$.
		\item \textbf{Grand Canonical (Macrocanonical):} $(\langle E \rangle, V, \langle N \rangle)$ fixed. Thermal and particle contact $(\mu, T)$.
	\end{itemize}

	\printbibliography[title={Bibliografía}]
\end{document}